\subsubsection{StandardPage}
\label{sec:StandardPage}
\paragraph{Descrizione}
La \textbf{StandardPage} è la pagina che funge da struttura del layout che ogni pagina riceve come parametro con il nome del componente \textbf{BasePageComponent}.
\paragraph{SottoComponenti}
StandardPage è composto dai seguenti sotto-componenti:
\begin{itemize}
    \item \textbf{BasePageComponent}: Il componente con dichiarati i metodi e parametri necessari alla creazione di una \textbf{StandardPage} tra i quali:
    \begin{itemize}
        \item \textbf{setHeader()}: Metodo che va a chiamare i componenti dediti alla costruzione del header;
        \item \textbf{setFooter()}: Metodo che va a chiamare i componenti dediti alla costruzione del footer;
        \item \textbf{setMessages()}: Metodo che richiama i componenti dedicati alla costruzione di messaggi di sistema, come messaggi a comparsa, pop-up e messaggi d'errore standardizzabili;
        \item \textbf{setLoadingStatus()}: Metodo che va a chiamare i componenti dediti al feedback del caricamento a livello grafico, ogni pagina dunque sarà capace di avviare un caricamento.
    \end{itemize}
    \item \textbf{ObserverComponent}: Componente che implementa il metodo \textbf{Update()}, chiamato dal \textbf{SubjectComponent} per reagire a cambiamenti e aggiornare il proprio stato o contenuto;
    \item \textbf{SubjectComponent}: Componente che gestisce un insieme di osservatori e offre il metodo \textbf{Subscribe()} per permettere loro di registrarsi e ricevere aggiornamenti;
    \item \textbf{NavigationLinksComponent}: Componente dotato del metodo \textbf{Update()} che eredita da ObserverComponent, tale componente si occupa di aggiornare i link presenti nel header a seconda della pagina richiesta;
    \item \textbf{HeaderComponent}: Componente dotato del metodo \textbf{Subscribe()} che eredita da SubjectComponent, tale componente riceve le modifiche richieste dal NavigationLinksComponent per aggiornare l'header;
    \item \textbf{FooterComponent}: Componente che istanzia il footer con la chiamata del metodo \textbf{ngOnInit()};
    \item \textbf{LoadingComponent}: Componente che dichiara i metodi necessari per l'inizializzazione del caricamento grafico:
    \begin{itemize}
        \item \textbf{loadingStatus}: Metodo booleano che indica la presenza o meno del caricamento;
        \item \textbf{ngOnInit()}: Metodo che si occupa di impostare il parametro \textbf{loadingStatus} a False;
        \item \textbf{setStatus()}: Metodo che imposta il valore del parametro \textbf{loadingStatus}.
    \end{itemize}
    \item \textbf{MessageHandlerComponent}: Componente dotato del metodo \textbf{Update()} che eredita da \textbf{ObserverComponent}, tale componente si occupa di creare il tipo e contenuto del messaggio;
    \item \textbf{MessageComponent}: Componente che eredita da \textbf{SubjectComponent} e implementa il metodo \textbf{Subscribe()}. Questo componente riceve le impostazioni del messaggio richieste dal \textbf{MessageHandlerComponent} per aggiornare la visualizzazione del messaggio.
\end{itemize}
